% Estrutura inicial do artigo do projeto diffusion_detection.
% Compilacao recomendada: pdflatex artigo_estrutura_inicial.tex (duas vezes).
% Antes da submissao, adaptar o arquivo ao template oficial do evento ou periodico.

\documentclass[11pt,a4paper]{article}

\usepackage[utf8]{inputenc}
\usepackage[T1]{fontenc}
\usepackage[brazil]{babel}
\usepackage{lmodern}
\usepackage{microtype}
\usepackage{geometry}
\usepackage{amsmath,amssymb}
\usepackage{booktabs}
\usepackage{array}
\usepackage{enumitem}
\usepackage{xcolor}
\usepackage{graphicx}
\usepackage{tikz}
\usetikzlibrary{arrows.meta,calc,fit,positioning}
\usepackage{hyperref}
\usepackage{url}

\geometry{margin=2.5cm}
\hypersetup{
  colorlinks=true,
  linkcolor=blue!45!black,
  citecolor=green!35!black,
  urlcolor=blue!55!black
}
\setlist{itemsep=0.25em, topsep=0.4em}
\renewcommand{\arraystretch}{1.2}

% Marcadores de trabalho para a equipe. Remover antes da submissao.
\newcommand{\afazer}[1]{%
  \par\smallskip
  \noindent\textcolor{red!70!black}{\textbf{A fazer pela equipe:} #1}
  \par\smallskip
}
\newcommand{\decidir}[1]{%
  \par\smallskip
  \noindent\textcolor{orange!70!black}{\textbf{Decisão pendente:} #1}
  \par\smallskip
}
\newcommand{\naoinventar}{%
  \par\smallskip
  \noindent\fbox{%
    \parbox{0.94\linewidth}{%
      \textbf{Nota de integridade científica:} não inserir métricas, resultados,
      comparações ou conclusões que ainda não tenham sido obtidos por uma
      execução versionada e auditável do experimento.}}
  \par\smallskip
}

\title{%
  Sinais Forenses Interpretáveis para Detecção de Vídeos Gerados por
  Inteligência Artificial sob Mudança de Gerador\\[0.4em]
  \large Estrutura inicial do manuscrito e plano de trabalho
}

\author{%
  Nome do(a) orientador(a)\\
  \texttt{email@instituicao.br}
  \and
  Nome dos(as) discentes\\
  Instituição
}

\date{Versão de trabalho --- setembro de 2026}

\begin{document}

\maketitle

\begin{abstract}
Os avanços recentes em modelos de geração de vídeo tornaram a distinção entre
conteúdo real e sintético progressivamente mais difícil, especialmente quando
o detector é avaliado em geradores não observados durante o treinamento. Este
trabalho investiga uma representação forense explícita e auditável, composta
por sinais de textura, bordas, estrutura local, resíduo, frequência,
fotometria e comportamento temporal, extraídos de regiões como rosto, olhos,
boca, corpo e fundo. A hipótese central é que a combinação de pistas de baixo
nível pode capturar inconsistências compartilhadas por diferentes geradores,
desde que a avaliação controle atalhos relacionados a conteúdo, identidade,
movimento, taxa de quadros, resolução e compressão. O estudo será conduzido
principalmente no benchmark DF26, usando partições agrupadas por conteúdo e um
protocolo de gerador não visto. \textbf{[Inserir aqui, somente após os
experimentos, os resultados principais com métricas e intervalos de
confiança.]}

\smallskip
\noindent\textbf{Palavras-chave:} vídeo sintético; detecção de conteúdo gerado
por IA; forense multimídia; sinais interpretáveis; generalização; mudança de
domínio.
\end{abstract}

\naoinventar

\section{Introdução}

Modelos contemporâneos de texto-para-vídeo e imagem-para-vídeo produzem cenas
com alto grau de realismo visual e temporal. Esse avanço reduz a utilidade de
detectores desenvolvidos apenas para manipulações faciais tradicionais, como
troca de rosto e reenactment. O problema torna-se particularmente relevante em
vídeos de pronunciamentos, entrevistas e comunicação pública, nos quais a
autenticidade do conteúdo pode influenciar decisões sociais e políticas.

Um obstáculo central é a generalização. Um detector pode apresentar desempenho
elevado quando os dados de treino e teste compartilham geradores, fontes,
codecs ou padrões de movimento, mas falhar diante de um novo processo
generativo. O benchmark DF26 evidencia essa dificuldade ao reunir vídeos de
fala produzidos por geradores recentes e ao reportar desempenho próximo ao
acaso para seres humanos e diferentes detectores existentes
\cite{shykula2026df26}. Estudos recentes também mostram que parte do desempenho
de detectores temporais pode decorrer de vieses de amostragem, movimento,
resolução ou pré-processamento, e não necessariamente de propriedades
intrínsecas da geração sintética \cite{michels2026biases}.

Este trabalho parte de uma direção complementar aos modelos profundos de ponta
a ponta. Em vez de aprender toda a representação diretamente dos pixels, o
pipeline calcula descritores determinísticos e semanticamente nomeados. Esses
descritores são agregados por vídeo e fornecidos a classificadores tabulares.
Essa separação permite analisar quais famílias de sinais e quais regiões
contribuem para a decisão, além de facilitar testes de robustez e a identificação
de atalhos.

A pergunta central do artigo é: \emph{até que ponto sinais forenses explícitos,
regionais e temporais discriminam vídeos reais e sintéticos quando o gerador de
teste não aparece no treinamento e os principais fatores de confusão são
controlados?}

As contribuições pretendidas são:

\begin{enumerate}
  \item uma representação tabular forense organizada por família de sinal,
        região semântica e comportamento temporal;
  \item um protocolo de avaliação que combine mudança de gerador e separação
        de conteúdo relacionado;
  \item uma análise de ablação por família, região e temporalidade;
  \item uma auditoria de atalhos associados a metadados, movimento,
        amostragem, resolução, codec e compressão.
\end{enumerate}

\afazer{Reescrever o último parágrafo das contribuições após os experimentos.
Cada contribuição declarada deverá apontar para uma tabela, figura, artefato de
código ou resultado verificável.}

\section{Fundamentação teórica}

\subsection{Detecção de vídeos totalmente gerados}

A literatura de deepfakes foi inicialmente dominada por manipulações locais do
rosto. Vídeos totalmente gerados constituem um regime diferente: aparência,
fundo, iluminação e movimento podem ser sintetizados conjuntamente. Por isso,
artefatos de recorte ou mistura de face, comuns em bases anteriores, não são
necessariamente informativos nesse novo cenário.

Abordagens recentes utilizam redes espaço-temporais, representações aprendidas,
modelos de estado e detectores baseados em frequência. O DeMamba, por exemplo,
explora inconsistências espaciais e temporais em um benchmark de grande escala
\cite{chen2024demamba}. O UNITE procura cobrir tanto manipulações parciais
quanto vídeos totalmente sintéticos \cite{kundu2025unite}. Já estratégias de
orientação forense procuram conduzir o modelo a pistas de baixo nível que
sejam compartilhadas por arquiteturas generativas distintas
\cite{corvi2025seeing}.

Esses resultados não eliminam o problema de mudança de domínio. Geradores
novos, recompressão de plataformas e diferenças entre as fontes reais e
sintéticas podem alterar as distribuições observadas. Consequentemente, o
desempenho dentro de uma mesma base não é evidência suficiente de uma propriedade
universal dos vídeos gerados.

\afazer{Construir uma tabela de trabalhos relacionados contendo: método,
representação, datasets, tipo de split, geradores não vistos, controles de
compressão/movimento, métricas, disponibilidade de código e limitações. Consultar
os artigos completos, não apenas resumos ou páginas secundárias.}

\subsection{Forense interpretável e sinais de baixo nível}

Uma representação interpretável não significa que cada atributo possua poder
causal ou seja isoladamente suficiente para detectar conteúdo sintético.
Significa que a grandeza calculada tem definição conhecida, unidade de análise
explícita e limites de interpretação documentados. Neste projeto, os sinais são
organizados em cinco grupos.

\subsubsection{Grupo A: textura, bordas e nitidez}

O Local Binary Pattern (LBP) descreve a relação entre um pixel e sua vizinhança
e é uma referência clássica para caracterização de textura
\cite{ojala2002lbp}. Para uma vizinhança de $P$ pontos e raio $R$, uma forma
geral do código é

\begin{equation}
  \operatorname{LBP}_{P,R}
  = \sum_{p=0}^{P-1} s(g_p-g_c)2^p,
  \qquad
  s(x)=
  \begin{cases}
    1, & x\geq 0,\\
    0, & x<0.
  \end{cases}
\end{equation}

O pipeline utiliza escalas distintas e resume propriedades das distribuições
dos códigos, como entropia normalizada, uniformidade e ocupação dos bins. Como
os códigos LBP são categorias, momentos numéricos arbitrários dos próprios
códigos não devem receber interpretação física.

Gradientes de Sobel representam intensidade e orientação de transições locais.
O Laplaciano responde à segunda derivada espacial e fornece indicadores de
bordas finas e nitidez. Essas grandezas também são sensíveis a blur,
sharpening, resize e compressão; portanto, diferenças entre classes precisam
ser reavaliadas sob perturbações controladas.

\subsubsection{Grupo B: estrutura local}

O SIFT detecta pontos de interesse e descreve padrões locais com relativa
robustez a escala e rotação \cite{lowe2004sift}. O projeto resume quantidade,
densidade, cobertura espacial, resposta, escala, orientação e autossimilaridade
dos descritores em cada região. Esses valores podem refletir estrutura local,
mas também conteúdo, pose, dimensão da região e qualidade de imagem.

A similaridade entre patches complementa o SIFT ao medir repetição de padrões
locais. A amostragem deve ser determinística e espacialmente distribuída, e as
contagens de patches válidos devem ser tratadas como controle de qualidade.

\subsubsection{Grupo C: resíduo de alta passagem}

O resíduo bilateral é definido por

\begin{equation}
  R = I - \operatorname{Bilateral}(I),
\end{equation}

em que o filtro bilateral suaviza variações preservando parte das bordas
\cite{tomasi1998bilateral}. Estatísticas de $R$ descrevem magnitude, dispersão,
caudas, autocorrelação espacial e correlação entre canais.

Esse resíduo contém uma mistura de textura, bordas, ruído, sharpening e
compressão. Ele deve ser apresentado como um \emph{baseline} de alta passagem,
e não como estimativa de PRNU, Noiseprint ou ruído de sensor.

\subsubsection{Grupo D: frequência espacial}

Para uma região em tons de cinza, o pipeline centraliza o sinal, aplica uma
janela de Hann e calcula a transformada discreta de Fourier:

\begin{equation}
  \widetilde{I}(x,y)=(I(x,y)-\overline{I})w(x,y),
  \qquad
  F(u,v)=\mathcal{F}_{2}\{\widetilde{I}\},
  \qquad
  P(u,v)=|F(u,v)|^2.
\end{equation}

Razões de potência por banda, centroide radial, entropia, planicidade,
anisotropia e inclinação espectral resumem a distribuição de energia. Pistas
frequenciais podem revelar artefatos de arquiteturas generativas, mas também
são fortemente afetadas por codec, taxa de compressão e interpolação. Trabalhos
recentes mostram tanto o potencial das pistas de frequência quanto a necessidade
de testá-las após recompressão \cite{corvi2025seeing,michels2026biases}.

\subsubsection{Grupo E: fotometria regional}

As métricas fotométricas resumem luminância, cromaticidade, saturação,
iluminação suavizada, gradientes e proporções de pixels escuros, claros ou
clipados. Assimetrias horizontais e verticais permitem descrever desequilíbrios
entre partes de uma região.

O módulo de candidatos de sombra usa uma aproximação inspirada em Retinex para
identificar pixels abaixo de uma estimativa local de iluminação. Esses valores
são \emph{proxies} fotométricos: não estimam completamente fonte de luz,
normais, reflectância ou geometria tridimensional. Materiais escuros e exposição
local podem gerar respostas semelhantes.

\subsection{Regiões semânticas e contrastes regionais}

Os descritores são extraídos de rosto completo, olhos, boca, corpo e fundo. A
hipótese é que uma inconsistência possa ser localizada ou aparecer como
contraste entre regiões. Entretanto, regiões incorretas, faces pequenas,
oclusões e fallbacks geométricos podem produzir separação artificial.

Cada observação deve preservar vídeo, frame, tempo, identidade rastreada,
região, origem da detecção e indicadores de qualidade. O MediaPipe fornece a
base operacional para detecção e segmentação \cite{lugaresi2019mediapipe}, mas
a validade científica das regiões depende de auditoria visual e quantitativa.

\subsection{Temporalidade}

Se $x(t)$ representa um sinal regional amostrado no tempo físico, as primeiras
e segundas diferenças podem ser aproximadas por

\begin{equation}
  \Delta x_i = \frac{x(t_{i+1})-x(t_i)}{t_{i+1}-t_i},
  \qquad
  \Delta^2 x_i =
  \frac{\Delta x_{i+1}-\Delta x_i}{\widetilde{t}_{i+1}-\widetilde{t}_i}.
\end{equation}

O pipeline também calcula dispersão robusta, autocorrelação e estatísticas de
frequência temporal. Usar timestamps é fundamental porque índices de frames
podem representar intervalos físicos diferentes. Ainda assim, menor ou maior
variação temporal pode ser apenas consequência de movimento, duração ou FPS.
Por isso, a contribuição temporal deverá ser avaliada antes e depois do
balanceamento desses fatores.

\subsection{Generalização, atalhos e vazamento experimental}

Atalhos são variáveis que permitem prever o rótulo dentro de uma base sem
representar a propriedade científica de interesse. Exemplos incluem resolução,
codec, FPS, duração, origem, quantidade de movimento, padding e política de
amostragem. Detectores temporais recentes mostraram quedas substanciais quando
vieses de movimento e pré-processamento foram controlados
\cite{michels2026biases}.

Outro risco é o vazamento de conteúdo. Se versões reais e sintéticas derivadas
do mesmo clipe, prompt ou identidade aparecem em treino e teste, o modelo pode
reconhecer o conteúdo em vez do processo de geração. Por isso, o artigo propõe
separar simultaneamente famílias de geradores e grupos de conteúdo.

\subsection{Classificadores tabulares e interpretação}

Após a extração, cada vídeo é representado por um vetor $z_V$, e um
classificador estima

\begin{equation}
  \widehat{y}=g_{\theta}(z_V).
\end{equation}

Regressão logística regularizada, Random Forest e gradient boosting representam
baselines com diferentes relações entre complexidade e interpretação
\cite{breiman2001rf,chen2016xgboost}. Imputação, normalização, seleção de
atributos e calibração devem ser ajustadas exclusivamente nos dados de treino
de cada fold. Coeficientes padronizados e importância por permutação agrupada
podem apoiar a análise, mas não demonstram causalidade.

\section{Perguntas de pesquisa e hipóteses}

\begin{description}[style=nextline]
  \item[RQ1] Sinais forenses explícitos discriminam vídeos reais e sintéticos
  quando o gerador de teste não aparece no treinamento?

  \item[RQ2] Quais famílias de sinais permanecem estáveis entre geradores?

  \item[RQ3] Quais regiões oferecem evidência complementar e quais capturam
  predominantemente conteúdo ou erro de segmentação?

  \item[RQ4] Os derivados temporais acrescentam informação após controlar FPS,
  duração e movimento?

  \item[RQ5] O desempenho persiste sob compressão, resize, mudanças de FPS e
  separação de conteúdo relacionado?
\end{description}

As hipóteses iniciais são:

\begin{enumerate}[label=H\arabic*.]
  \item A representação forense supera um baseline baseado apenas em metadados
        nos testes com gerador não visto.
  \item Sinais espectrais e residuais apresentam maior transferência entre
        geradores, mas são sensíveis a recompressão e resize.
  \item A temporalidade oferece ganho incremental apenas quando amostragem e
        movimento são controlados.
  \item Contrastes regionais acrescentam informação sem depender
        predominantemente de fallback ou falha do detector de regiões.
\end{enumerate}

\afazer{Congelar perguntas e hipóteses antes da análise confirmatória. Se uma
hipótese for refutada, manter e relatar a refutação; não reescrever a hipótese
para coincidir com os resultados.}

\section{Materiais e métodos}

A Figura~\ref{fig:visao-geral-metodo} sintetiza o método proposto e explicita
a separação entre a representação forense, composta por operações
determinísticas, e a etapa supervisionada. Variáveis de aquisição e governança
não compõem $z_V$; elas são mantidas para auditoria e para o baseline de
atalhos. Essa separação permite atribuir eventuais ganhos às famílias de sinais
e verificar se o desempenho persiste sob mudança de gerador e de conteúdo.

\begin{figure}[htbp]
  \centering
  \resizebox{\linewidth}{!}{% Figura vetorial da arquitetura metodologica.
% O arquivo e incluído por artigo_estrutura_inicial.tex dentro de um resizebox.
\begin{tikzpicture}[
  font=\sffamily\scriptsize,
  node distance=4.5mm and 5.5mm,
  >=Latex,
  fluxo/.style={->, line width=0.75pt, draw=black!65},
  fluxoaux/.style={->, line width=0.65pt, draw=black!45, dashed},
  bloco/.style={
    draw=black!55,
    line width=0.65pt,
    rounded corners=2pt,
    align=center,
    inner sep=4.5pt,
    minimum height=10mm
  },
  entrada/.style={bloco, fill=blue!8},
  processo/.style={bloco, fill=black!3},
  aprendido/.style={bloco, fill=orange!13, draw=orange!65!black},
  saida/.style={bloco, fill=green!11, draw=green!55!black},
  familia/.style={
    bloco,
    text width=26mm,
    minimum height=15mm,
    inner sep=3.5pt
  },
  rotulo/.style={font=\sffamily\bfseries\scriptsize, text=black!65},
  nota/.style={
    draw=black!38,
    dashed,
    rounded corners=2pt,
    fill=black!1,
    align=left,
    inner sep=4.5pt
  }
]

% Entrada e decomposicao do video.
\node[entrada, text width=25mm] (video) {\textbf{Vídeo de entrada}\\$V=\{I_1,\ldots,I_T\}$};
\node[processo, text width=37mm, right=of video] (amostragem) {
  \textbf{Amostragem em tempo físico}\\
  timestamps, padronização e controle de qualidade
};
\node[processo, text width=44mm, right=of amostragem] (regioes) {
  \textbf{Regiões semânticas}\\
  rosto, olhos, boca, corpo e fundo\\
  \textcolor{black!62}{track, confiança, origem e fallback}
};

\draw[fluxo] (video) -- (amostragem);
\draw[fluxo] (amostragem) -- (regioes);

% Banco de descritores determinísticos.
\node[familia, fill=blue!9, below=12mm of video, xshift=-2mm] (ga) {
  \textbf{A · Textura e bordas}\\
  LBP, Sobel e Laplaciano
};
\node[familia, fill=cyan!10, right=3.2mm of ga] (gb) {
  \textbf{B · Estrutura local}\\
  SIFT e similaridade de patches
};
\node[familia, fill=violet!9, right=3.2mm of gb] (gc) {
  \textbf{C · Resíduo}\\
  alta passagem bilateral
};
\node[familia, fill=orange!11, right=3.2mm of gc] (gd) {
  \textbf{D · Frequência}\\
  potência e geometria espectral
};
\node[familia, fill=green!10, right=3.2mm of gd] (ge) {
  \textbf{E · Fotometria}\\
  luminância, cor e iluminação
};

\coordinate (familias-top) at ($(gc.north)+(0,4.5mm)$);
\draw[fluxo] (regioes.south) |- (familias-top);
\draw[line width=0.75pt, draw=black!65] (ga.north |- familias-top) -- (ge.north |- familias-top);
\foreach \n in {ga,gb,gc,gd,ge}{
  \draw[fluxo] (\n.north |- familias-top) -- (\n.north);
}

\node[processo, text width=44mm, below=11mm of gb, xshift=-5mm] (sinais) {
  \textbf{Sinais regionais por quadro}\\
  $x_{V,t,r,k}=\phi_k(I_{V,t},r)$\\
  \textcolor{black!62}{inclui contrastes entre regiões}
};
\node[processo, text width=54mm, right=7mm of sinais] (agregacao) {
  \textbf{Agregação estática e temporal}\\
  média, desvio e mediana; $\Delta x$, $\Delta^2x$,\\
  autocorrelação e espectro temporal
};
\node[entrada, text width=24mm, right=7mm of agregacao] (vetor) {
  \textbf{Vetor por vídeo}\\
  $z_V=\Psi(\Phi(I_{1:T}))$
};

\coordinate (familias-bottom) at ($(gc.south)+(0,-4.5mm)$);
\draw[line width=0.75pt, draw=black!65] (ga.south |- familias-bottom) -- (ge.south |- familias-bottom);
\foreach \n in {ga,gb,gc,gd,ge}{
  \draw[fluxo] (\n.south) -- (\n.south |- familias-bottom);
}
\draw[fluxo] (familias-bottom) -| (sinais.north);
\draw[fluxo] (sinais) -- (agregacao);
\draw[fluxo] (agregacao) -- (vetor);

% Etapa aprendida. O fluxo retorna da direita para a esquerda para evitar
% cruzamentos entre as duas faixas da figura.
\node[aprendido, text width=39mm, below=13mm of vetor] (ajuste) {
  \textbf{Ajuste dentro de cada fold}\\
  imputação, escala, seleção e calibração
};
\node[aprendido, text width=43mm, left=7mm of ajuste] (classificador) {
  \textbf{Classificador tabular $g_\theta$}\\
  regressão logística, Random Forest\\
  ou gradient boosting
};
\node[saida, text width=28mm, left=7mm of classificador] (predicao) {
  \textbf{Predição}\\
  $\widehat p(y=\mathrm{fake}\mid V)$
};

\draw[fluxo] (vetor) -- (ajuste);
\draw[fluxo] (ajuste) -- (classificador);
\draw[fluxo] (classificador) -- (predicao);

% Protocolo e auditoria: influencia a avaliacao, mas nao e entrada forense.
\node[nota, text width=137mm, below=11mm of classificador] (protocolo) {
  \textbf{Protocolo de avaliação e auditoria (fora de $z_V$):}
  separação por conteúdo + gerador não visto (LOGOS); baseline de metadados;
  ablações por família, região e temporalidade; controles de movimento, FPS,
  resolução, codec e compressão; holdout comercial após congelamento.
};

\draw[fluxoaux] (protocolo.north -| ajuste.south) -- (ajuste.south);
\draw[fluxoaux] (protocolo.north -| predicao.south) -- (predicao.south);

% Contornos que explicitam o que e fixo e o que e aprendido.
\node[
  draw=blue!55!black,
  line width=0.75pt,
  dashed,
  rounded corners=3pt,
  fit=(amostragem)(regioes)(ga)(ge)(sinais)(agregacao)(vetor),
  inner sep=4.5mm,
  label={[rotulo, text=blue!55!black]above left:extração determinística e auditável}
] (deterministico) {};

\node[
  draw=orange!70!black,
  line width=0.75pt,
  dashed,
  rounded corners=3pt,
  fit=(ajuste)(classificador)(predicao),
  inner sep=4.5mm,
  label={[rotulo, text=orange!70!black]above left:etapa supervisionada}
] (supervisionado) {};

\end{tikzpicture}
}
  \caption{Visão geral do método. Cada vídeo é amostrado em tempo físico e
  decomposto em regiões semânticas. Cinco famílias de descritores produzem
  sinais por quadro e região, posteriormente resumidos por estatísticas
  estáticas e temporais em um vetor por vídeo $z_V$. Apenas a etapa final é
  aprendida. O protocolo de avaliação e as variáveis de auditoria não são
  fornecidos ao detector forense.}
  \label{fig:visao-geral-metodo}
\end{figure}

\subsection{Estado atual do repositório}

O repositório implementa ingestão, pré-processamento regional, extração dos
grupos A--E, agregações temporais, camadas Bronze/Silver/Gold, contratos de
dados e versionamento por DVC. No momento de preparação desta estrutura, não
estavam disponíveis localmente a base DF26 processada, a Gold experimental, um
módulo completo de treinamento em \texttt{src/ml} ou resultados persistidos.
Os notebooks também não continham saídas executadas.

Portanto, as seções de resultados e conclusão deste documento são apenas uma
estrutura de preenchimento. O artefato registrado no \texttt{dvc.lock}
corresponde a validação operacional e não deve ser tratado como amostra
científica.

\subsection{Dados}

O DF26 será a base principal. O benchmark contém 271 vídeos reais e 2.420
vídeos sintéticos de sete famílias de geradores em cenários de fala com uma
pessoa \cite{shykula2026df26}. A licença e os papéis experimentais definidos
pela base devem ser respeitados, mantendo geradores comerciais apenas para
avaliação final.

O catálogo adicional do projeto contém 2.039 links, dos quais 1.072 estão
rotulados como reais e 967 como falsos. Essa coleção somente poderá ser usada
como validação externa após auditoria de disponibilidade, licença, proveniência,
rótulos, duplicatas e vídeos derivados.

\afazer{Produzir a tabela final de composição da amostra após todos os filtros,
com contagens por classe, gerador, cenário, modo de geração, resolução, FPS,
codec e motivo de exclusão.}

\subsection{Pré-processamento e regiões}

Os vídeos são amostrados em tempo físico. O pipeline atual prevê 5 fps e até 25
frames no DF26, preservando FPS, duração, contagem de frames e resolução
originais para auditoria. MediaPipe é usado para produzir regiões de rosto,
olhos, boca, corpo e fundo, com registro de confiança, origem da região e
fallback.

\afazer{Validar a política de amostragem, auditar visualmente uma amostra
estratificada e definir critérios objetivos para face pequena, oclusão, pose,
blur, troca de identidade e fallback.}

\subsection{Extração e agregação dos sinais}

Para o vídeo $V$, frame $t$, região $r$ e descritor $k$, define-se

\begin{equation}
  x_{V,t,r,k}=\phi_k(I_{V,t},r).
\end{equation}

As observações são agregadas por vídeo, região e track, gerando estatísticas
estáticas e temporais:

\begin{equation}
  z_V=\Psi\bigl(\Phi(I_1),\Phi(I_2),\ldots,\Phi(I_T)\bigr).
\end{equation}

Colunas de governança, caminhos, nomes de arquivo, fonte e identidade do
gerador não podem entrar como atributos preditivos. Variáveis de aquisição
serão reservadas para auditoria e para um baseline de atalhos.

\subsection{Protocolo de avaliação}

O protocolo LOGOS atualmente implementado precisa ser corrigido: os folds
existentes colocam todos os vídeos reais em treino e somente falsos do gerador
omitido em teste. Isso não permite calcular uma métrica binária real-versus-falso
válida no conjunto de teste.

O protocolo proposto terá dois eixos:

\begin{enumerate}
  \item \textbf{mudança de gerador:} uma família open-weight é omitida do
        treinamento;
  \item \textbf{mudança de conteúdo:} \texttt{df26\_clip\_id} e demais chaves
        de identidade/conteúdo são divididas em grupos disjuntos.
\end{enumerate}

Cada teste deverá conter vídeos reais e falsos. Treino, validação e teste não
podem compartilhar clipes relacionados. O conjunto comercial será aberto uma
única vez, depois que representação, modelos, hiperparâmetros e limiares forem
congelados.

\afazer{Implementar o novo gerador de splits, adicionar testes automatizados de
ausência de sobreposição e salvar a tabela de grupos usada em cada repetição.}

\subsection{Baselines e ablações}

Os experimentos mínimos serão:

\begin{table}[ht]
  \centering
  \caption{Experimentos mínimos planejados.}
  \label{tab:experimentos}
  \begin{tabular}{p{0.28\linewidth}p{0.62\linewidth}}
    \toprule
    Experimento & Finalidade \\
    \midrule
    Aleatório e majoritário & Referências inferiores de desempenho. \\
    Metadados somente & Quantificar atalhos de aquisição. \\
    Regressão logística & Baseline linear regularizado. \\
    Random Forest/boosting & Baselines tabulares não lineares. \\
    Grupos A--E isolados & Medir contribuição de cada família. \\
    Estático vs. temporal & Medir ganho temporal incremental. \\
    Regiões isoladas & Localizar a fonte da evidência. \\
    Sem fallback & Testar dependência de falhas regionais. \\
    Perturbações & Medir robustez a compressão, resize, FPS, blur e ruído. \\
    Holdout comercial & Avaliação final com pipeline congelado. \\
    \bottomrule
  \end{tabular}
\end{table}

\subsection{Métricas e análise estatística}

Serão reportadas AUROC, macro-AUROC entre geradores, balanced accuracy, AUPRC
com prevalência, TPR em FPR fixo, Brier score e erro esperado de calibração.
Também serão medidos tempo de extração, memória e dimensionalidade.

Intervalos de confiança serão estimados por bootstrap agrupado em
\texttt{df26\_clip\_id}. Comparações de modelos usarão as mesmas unidades de
teste. Análises univariadas deverão corrigir múltiplas comparações. O vídeo é a
unidade mínima de avaliação; frames do mesmo vídeo não serão tratados como
amostras independentes.

\afazer{Definir previamente número de repetições, sementes, FPR operacional,
método de correção de múltiplos testes e regra de seleção do modelo final.}

\section{Plano de trabalho até o artigo}

\begin{table}[ht]
  \centering
  \small
  \caption{Pacotes de trabalho e evidências de conclusão.}
  \label{tab:plano}
  \begin{tabular}{p{0.05\linewidth}p{0.21\linewidth}p{0.45\linewidth}p{0.17\linewidth}}
    \toprule
    ID & Frente & Entrega verificável & Responsável \\
    \midrule
    P1 & Protocolo & Splits agrupados, testes de vazamento e plano estatístico congelado. & [nome] \\
    P2 & Ambiente & Dependências fixadas; testes e pipeline reproduzidos em ambiente limpo. & [nome] \\
    P3 & Dados & DF26 obtido sob licença; manifestos, hashes e exclusões registrados. & [nome] \\
    P4 & Regiões & Auditoria visual estratificada e relatório de cobertura/fallback. & [nome] \\
    P5 & EDA & Relatório de qualidade, confundidores, correlação e estabilidade. & [nome] \\
    P6 & Modelagem & Baselines, seleção interna aos folds e calibração implementados. & [nome] \\
    P7 & Ablacões & Resultados por família, região, temporalidade e gerador. & [nome] \\
    P8 & Robustez & Compressão, resize, FPS, movimento e controles negativos. & [nome] \\
    P9 & Comparação & Pelo menos um método externo sob protocolo compatível. & [nome] \\
    P10 & Redação & Figuras, tabelas, discussão, limitações e pacote reproduzível. & [nome] \\
    \bottomrule
  \end{tabular}
\end{table}

\subsection{Ordem recomendada}

\begin{enumerate}
  \item corrigir o protocolo antes de consultar resultados confirmatórios;
  \item estabilizar o ambiente e executar os testes;
  \item processar um piloto estratificado e auditar as regiões;
  \item executar EDA de qualidade e corrigir problemas de medição;
  \item processar a base completa e treinar os baselines;
  \item executar ablações e controles de robustez;
  \item congelar o pipeline e avaliar o holdout comercial;
  \item preencher resultados, discussão, limitações e conclusão.
\end{enumerate}

\subsection{Critérios mínimos para submissão}

O manuscrito somente deverá ser submetido quando:

\begin{itemize}
  \item cada conjunto de teste contiver ambas as classes;
  \item não houver sobreposição indevida de conteúdo, identidade ou derivados;
  \item a cobertura e as falhas das regiões estiverem documentadas;
  \item os experimentos forem reproduzíveis em ambiente limpo;
  \item o ganho sobre baselines simples e metadados persistir em geradores não
        vistos, ou a ausência desse ganho estiver rigorosamente demonstrada;
  \item os principais confundidores tiverem sido controlados;
  \item o holdout comercial tiver sido usado apenas depois do congelamento;
  \item médias vierem acompanhadas de resultados por gerador e intervalos de
        confiança.
\end{itemize}

\section{Resultados}

\naoinventar

\subsection{Qualidade dos dados e das regiões}

\afazer{Inserir contagens, exclusões, cobertura regional, confiança, fallback,
tracks e falhas. Apresentar esses resultados antes das métricas preditivas.}

\subsection{Generalização para geradores não vistos}

\afazer{Inserir a tabela principal por fold e gerador, com intervalos de
confiança. Incluir o baseline de metadados na mesma comparação.}

\subsection{Ablacão de famílias e regiões}

\afazer{Relatar grupos A--E, regiões e contrastes. Discutir estabilidade entre
geradores, e não apenas o melhor valor médio.}

\subsection{Contribuição da temporalidade}

\afazer{Comparar modelos estáticos e temporais antes e depois de controlar FPS,
duração e movimento.}

\subsection{Robustez e avaliação comercial}

\afazer{Relatar degradação sob perturbações e, por último, o holdout comercial.
Registrar qualquer análise realizada após a abertura desse holdout.}

\section{Discussão}

A discussão deverá responder, com base nos resultados:

\begin{itemize}
  \item quais sinais permanecem úteis entre geradores;
  \item quais diferenças desaparecem quando confundidores são controlados;
  \item se regiões semânticas oferecem evidência complementar;
  \item se a temporalidade mede artefato generativo ou viés de movimento;
  \item qual é a relação entre desempenho, custo e auditabilidade;
  \item quais conclusões não podem ser generalizadas além do DF26.
\end{itemize}

Um resultado próximo ao acaso após controles não invalida automaticamente o
trabalho. Ele pode sustentar uma contribuição de auditoria ou resultado
negativo, desde que o protocolo seja sólido e a causa da queda seja analisada.

\afazer{Confrontar cada hipótese com os resultados e com trabalhos primários.
Separar evidência observada, interpretação plausível e especulação.}

\section{Limitações, ética e reprodutibilidade}

O DF26 é concentrado em vídeos curtos de uma pessoa falando, o que limita a
validade externa. O desempenho pode mudar em cenas longas, múltiplas pessoas,
câmeras móveis ou outros gêneros audiovisuais. Regiões incorretas, rótulos
imprecisos e versões futuras dos geradores constituem ameaças adicionais.

Detectores de conteúdo sintético têm implicações de uso dual. Uma previsão não
deve ser comunicada como prova definitiva de falsidade, especialmente em
contextos de alto impacto. O artigo deverá descrever incerteza, limitações e
condições de falha.

Para reprodutibilidade, serão publicados, dentro das restrições de licença,
código, parâmetros, versões, seeds, splits, hashes, critérios de exclusão e
resultados agregados. Os vídeos do DF26 não deverão ser redistribuídos fora dos
termos autorizados.

\section{Conclusão}

Este manuscrito propõe avaliar sinais forenses explícitos para detecção de
vídeos gerados por IA sob mudança de gerador e separação de conteúdo. A
estrutura atual do repositório permite conduzir essa investigação, mas ainda é
necessário corrigir o protocolo, processar a base experimental, auditar as
regiões, implementar a modelagem e executar ablações e testes de robustez.

\textbf{[Substituir este parágrafo após os experimentos por uma resposta direta
às perguntas de pesquisa, delimitada aos dados e controles efetivamente
avaliados.]}

\appendix

\section{Checklist da equipe antes de remover os marcadores}

\begin{itemize}
  \item[$\square$] Todas as referências foram conferidas nas fontes primárias.
  \item[$\square$] Toda afirmação quantitativa aponta para um artefato versionado.
  \item[$\square$] Nenhuma coluna de governança entrou como feature.
  \item[$\square$] Imputação, escala e seleção foram ajustadas dentro dos folds.
  \item[$\square$] O teste contém reais e falsos sem vazamento de conteúdo.
  \item[$\square$] Há resultados por gerador e intervalos de confiança.
  \item[$\square$] O baseline de metadados foi executado.
  \item[$\square$] As falhas de região e fallback foram auditadas.
  \item[$\square$] O holdout comercial foi aberto somente após o congelamento.
  \item[$\square$] Todas as ocorrências de ``A fazer'' foram resolvidas.
  \item[$\square$] O texto foi migrado para o template do evento ou periódico.
\end{itemize}

\begin{thebibliography}{99}

\bibitem{shykula2026df26}
S.~Shykula, A.~Yermakov, I.~Samarskyi, D.~Mishkin, J.~Cech e A.~Mishchuk.
\newblock DF26: We Cannot Tell Fake From Real Anymore.
\newblock \emph{arXiv preprint arXiv:2609.07369}, 2026.

\bibitem{michels2026biases}
J.~Michels, L.~Jorissen e N.~Michiels.
\newblock Dataset Biases and Shortcut Learning in Motion-Based AI-Generated
Video Detection.
\newblock \emph{arXiv preprint arXiv:2607.00948}, 2026.

\bibitem{corvi2025seeing}
R.~Corvi, D.~Cozzolino, E.~Prashnani, S.~De Mello, K.~Nagano e L.~Verdoliva.
\newblock Seeing What Matters: Generalizable AI-generated Video Detection with
Forensic-Oriented Augmentation.
\newblock \emph{arXiv preprint arXiv:2506.16802}, 2025.

\bibitem{chen2024demamba}
H.~Chen et al.
\newblock DeMamba: AI-Generated Video Detection on Million-Scale GenVideo
Benchmark.
\newblock \emph{arXiv preprint arXiv:2405.19707}, 2024.

\bibitem{kundu2025unite}
S.~Kundu et al.
\newblock Towards a Universal Synthetic Video Detector: From Face or
Background Manipulations to Fully AI-Generated Content.
\newblock In \emph{Proceedings of the IEEE/CVF Conference on Computer Vision
and Pattern Recognition}, 2025.

\bibitem{ojala2002lbp}
T.~Ojala, M.~Pietikäinen e T.~Mäenpää.
\newblock Multiresolution Gray-Scale and Rotation Invariant Texture
Classification with Local Binary Patterns.
\newblock \emph{IEEE Transactions on Pattern Analysis and Machine Intelligence},
24(7):971--987, 2002.

\bibitem{lowe2004sift}
D.~G. Lowe.
\newblock Distinctive Image Features from Scale-Invariant Keypoints.
\newblock \emph{International Journal of Computer Vision}, 60:91--110, 2004.

\bibitem{tomasi1998bilateral}
C.~Tomasi e R.~Manduchi.
\newblock Bilateral Filtering for Gray and Color Images.
\newblock In \emph{Proceedings of the Sixth International Conference on
Computer Vision}, pages 839--846, 1998.

\bibitem{lugaresi2019mediapipe}
C.~Lugaresi et al.
\newblock MediaPipe: A Framework for Building Perception Pipelines.
\newblock \emph{arXiv preprint arXiv:1906.08172}, 2019.

\bibitem{breiman2001rf}
L.~Breiman.
\newblock Random Forests.
\newblock \emph{Machine Learning}, 45:5--32, 2001.

\bibitem{chen2016xgboost}
T.~Chen e C.~Guestrin.
\newblock XGBoost: A Scalable Tree Boosting System.
\newblock In \emph{Proceedings of the 22nd ACM SIGKDD International Conference
on Knowledge Discovery and Data Mining}, pages 785--794, 2016.

\end{thebibliography}

\end{document}
